% GENERATED FILE. Edit canonical lessons, then run npm run generate:books.
% canonical-source-hash: fnv1a64:dde02fe276718aa5
% canonical-lessons: ES-C548-matrimonio, ES-C548-ceremonia, ES-C548-banquete, ES-C548-brindis, ES-C548-luna-de-miel

\chapter{El banquete}
\label{ch:es-el-banquete}
\hlchaptermodality{548}

\begin{chapteropening}
\textbf{By the end of this chapter:} \emph{I can describe a wedding: the marriage, the ceremony, the reception meal, the toast and the honeymoon.}

Complete the last lesson of chapter 548: I can describe a wedding: the marriage, the ceremony, the reception meal, the toast and the honeymoon.
\end{chapteropening}

\section[el matrimonio]{el matrimonio — the marriage; the married couple}

\label{lesson:ES-C548-matrimonio}



\begin{quote}
Before the new one: say the Spanish for the engagement, then the Spanish for the ring.
\end{quote}

\subsection*{You'll want to know: el matrimonio}
\textbf{el matrimonio} — \textquotedblleft{}the marriage; the married couple\textquotedblright{}.

\begin{quote}
Su matrimonio duró cincuenta años.
\end{quote}

\begin{quote}
\emph{Their marriage lasted fifty years.}
\end{quote}

\emph{Un matrimonio} is also a married couple: \emph{En el quinto vive un matrimonio mayor.}

\begin{cousinweb}
\emph{matrimonio} is a learned borrowing from Latin \emph{mātrimōnium} 'marriage, wedlock' (EtymDB), built on \emph{māter} 'mother'.
\end{cousinweb}

\begin{grammarlens}[title={What you've built}]
One more word for telling, explaining and describing at B1.
\end{grammarlens}

\subsection*{Your turn}
\begin{itemize}
\raggedright
  \item \emph{Say it:} \emph{el matrimonio}
  \item \emph{Say it:} \emph{el matrimonio}, once more
  \item \emph{Say it:} a sentence of your own with \emph{el matrimonio}
  \item \emph{Recall:} say the Spanish for the engagement, then the Spanish for the ring, then say \emph{el matrimonio} again
\end{itemize}

\begin{tcolorbox}[breakable,colback=teal!4,colframe=teal!35!black,title={Before you move on}]
What does el matrimonio mean? (\textquotedblleft{}The marriage; the married couple\textquotedblright{}.) Say it once more.
\end{tcolorbox}
\section[la ceremonia]{la ceremonia — the ceremony}

\label{lesson:ES-C548-ceremonia}



\begin{quote}
Before the new one: say the Spanish for the ring, then the Spanish for the marriage.
\end{quote}

\subsection*{You'll want to know: la ceremonia}
\textbf{la ceremonia} — \textquotedblleft{}the ceremony\textquotedblright{}.

\begin{quote}
La ceremonia fue corta y bonita.
\end{quote}

\begin{quote}
\emph{The ceremony was short and lovely.}
\end{quote}

\begin{cousinweb}
\emph{ceremonia} is a learned borrowing from Latin \emph{caerimōnia} 'a sacred rite' (EtymDB).
\end{cousinweb}

\begin{grammarlens}[title={What you've built}]
One more word for telling, explaining and describing at B1.
\end{grammarlens}

\subsection*{Your turn}
\begin{itemize}
\raggedright
  \item \emph{Say it:} \emph{la ceremonia}
  \item \emph{Say it:} \emph{la ceremonia}, once more
  \item \emph{Say it:} a sentence of your own with \emph{la ceremonia}
  \item \emph{Recall:} say the Spanish for the ring, then the Spanish for the marriage, then say \emph{la ceremonia} again
\end{itemize}

\begin{tcolorbox}[breakable,colback=teal!4,colframe=teal!35!black,title={Before you move on}]
What does la ceremonia mean? (\textquotedblleft{}The ceremony\textquotedblright{}.) Say it once more.
\end{tcolorbox}
\section[el banquete]{el banquete — the banquet, the wedding meal}

\label{lesson:ES-C548-banquete}



\begin{quote}
Before the new one: say the Spanish for the marriage, then the Spanish for the ceremony.
\end{quote}

\subsection*{You'll want to know: el banquete}
\textbf{el banquete} — \textquotedblleft{}the banquet, the wedding meal\textquotedblright{}.

\begin{quote}
Después de la boda hubo un banquete.
\end{quote}

\begin{quote}
\emph{After the wedding there was a banquet.}
\end{quote}

\begin{cousinweb}
\emph{banquete} is usually traced to French \emph{banquet} or Italian \emph{banchetto}, a small form of \emph{banco} 'bench', a Germanic word. The feast was named after the little benches the guests sat on.
\end{cousinweb}

\begin{grammarlens}[title={What you've built}]
One more word for telling, explaining and describing at B1.
\end{grammarlens}

\subsection*{Your turn}
\begin{itemize}
\raggedright
  \item \emph{Say it:} \emph{el banquete}
  \item \emph{Say it:} \emph{el banquete}, once more
  \item \emph{Say it:} a sentence of your own with \emph{el banquete}
  \item \emph{Recall:} say the Spanish for the marriage, then the Spanish for the ceremony, then say \emph{el banquete} again
\end{itemize}

\begin{tcolorbox}[breakable,colback=teal!4,colframe=teal!35!black,title={Before you move on}]
What does el banquete mean? (\textquotedblleft{}The banquet, the wedding meal\textquotedblright{}.) Say it once more.
\end{tcolorbox}
\section[el brindis]{el brindis — the toast (with a drink)}

\label{lesson:ES-C548-brindis}



\begin{quote}
Before the new one: say the Spanish for the ceremony, then the Spanish for the banquet.
\end{quote}

\subsection*{You'll want to know: el brindis}
\textbf{el brindis} — \textquotedblleft{}the toast (with a drink)\textquotedblright{}.

\begin{quote}
El padrino hizo un brindis por los novios.
\end{quote}

\begin{quote}
\emph{The best man made a toast to the couple.}
\end{quote}

The same form in the plural: \emph{los brindis}.

\begin{cousinweb}
\emph{brindis} comes from the German toast \emph{bring dir's} 'I bring it to you' (EtymDB), said to have been picked up from German soldiers in the 1500s.
\end{cousinweb}

\begin{grammarlens}[title={What you've built}]
One more word for telling, explaining and describing at B1.
\end{grammarlens}

\subsection*{Your turn}
\begin{itemize}
\raggedright
  \item \emph{Say it:} \emph{el brindis}
  \item \emph{Say it:} \emph{el brindis}, once more
  \item \emph{Say it:} a sentence of your own with \emph{el brindis}
  \item \emph{Recall:} say the Spanish for the ceremony, then the Spanish for the banquet, then say \emph{el brindis} again
\end{itemize}

\begin{tcolorbox}[breakable,colback=teal!4,colframe=teal!35!black,title={Before you move on}]
What does el brindis mean? (\textquotedblleft{}The toast (with a drink)\textquotedblright{}.) Say it once more.
\end{tcolorbox}
\section[la luna de miel]{la luna de miel — the honeymoon}

\label{lesson:ES-C548-luna-de-miel}



\begin{quote}
Before the new one: say the Spanish for the banquet, then the Spanish for the toast.
\end{quote}

\subsection*{You'll want to know: la luna de miel}
\textbf{la luna de miel} — \textquotedblleft{}the honeymoon\textquotedblright{}.

\begin{quote}
Pasaron la luna de miel en Cuba.
\end{quote}

\begin{quote}
\emph{They spent their honeymoon in Cuba.}
\end{quote}

\emph{Irse de luna de miel} — to go on honeymoon.

\begin{cousinweb}
\emph{luna de miel} was built on Latin \emph{lūna} 'moon' and \emph{mel} 'honey' (Lewis \& Short), and is usually said to copy English \emph{honeymoon}, through French \emph{lune de miel}.
\end{cousinweb}

\begin{grammarlens}[title={What you've built}]
One more word for telling, explaining and describing at B1.
\end{grammarlens}

\subsection*{Your turn}
\begin{itemize}
\raggedright
  \item \emph{Say it:} \emph{la luna de miel}
  \item \emph{Say it:} \emph{la luna de miel}, once more
  \item \emph{Say it:} a sentence of your own with \emph{la luna de miel}
  \item \emph{Recall:} say the Spanish for the banquet, then the Spanish for the toast, then say \emph{la luna de miel} again
\end{itemize}

\begin{tcolorbox}[breakable,colback=teal!4,colframe=teal!35!black,title={Before you move on}]
What does la luna de miel mean? (\textquotedblleft{}The honeymoon\textquotedblright{}.) Say it once more.
\end{tcolorbox}
